%%%%%%%%%%%%%%%%%%%%%%%%%%%%%%%%%%%%%%%%%%%%%%%%%%%%%%%%%%%%%%%%%%%%%%%
%%%%%%%%%%%%%%%%%%%%%%%%%%%%%%%%%%%%%%%%%%%%%%%%%%%%%%%%%%%%%%%%%%%%%%%
%% start of file `minimalthesis_logo.tex'.                           %%
%%                                                                   %%
%% version: v0.0.2-alpha3                                            %%
%% status: alpha                                                     %%
%% file version: 2023-12-18                                          %%
%%                                                                   %%
%% Copyright 2023 Daniel Sacré.                                      %%
%%                                                                   %%    
%% The official logo for the minimalthesis LaTeX Documentclass.      %%
%% natural sciences.                                                 %%
%%                                                                   %%
%% This piece of software comes without any warranty, not even that  %%
%% it is fit for the intended purpose. Since it is an alpha-release, %%
%% Bug reports and suggestions via github are highly welcome.        %%
%%                                                                   %%
%% LICENSE:                                                          %%
%% This work may be distributed and/or modified under the            %%
%% conditions of the Creative Commons Attribution, NonCommercial,    %%
%% NoDerivatives 4.0 International License                           %%
%%%%%%%%%%%%%%%%%%%%%%%%%%%%%%%%%%%%%%%%%%%%%%%%%%%%%%%%%%%%%%%%%%%%%%%
%%%%%%%%%%%%%%%%%%%%%%%%%%%%%%%%%%%%%%%%%%%%%%%%%%%%%%%%%%%%%%%%%%%%%%%

\documentclass{standalone}

\usepackage{tikz}
\usetikzlibrary{shadings,shapes,positioning,calc,patterns}
\usetikzlibrary{
	decorations.pathmorphing,
	decorations.pathreplacing,
	decorations.shapes,
	decorations.text,
}

% source: https://tex.stackexchange.com/questions/66216/draw-arc-in-tikz-when-center-of-circle-is-specified
\def\centerarc[#1](#2)(#3:#4:#5)% Syntax: [draw options] (center) (initial angle:final angle:radius)
{ \draw[#1] ($(#2)+({#5*cos(#3)},{#5*sin(#3)})$) arc (#3:#4:#5); }

\def\scale{0.55}

% source of typewriter image: https://openclipart.org/detail/322953/typewriter, License: CC0

\begin{document}
	\begin{tikzpicture}
		% old variant
%		\node[fill=cyan, circle, minimum width=8cm, text=white, align=center, font=\Huge]at(0,0){$<min>$\\Thesis};

		% new variant
		% outer background ring
		\path[fill=cyan, even odd rule](0,0)circle(\scale*5.7cm)(0,0)circle(\scale*4.125cm);
		
		% typewriter with filled inner background circle
		\node at (0,0){
			\includegraphics[width=4.4 cm]{typewriter_openclipart_edited.pdf}
		};
		
		% curved text top
		\centerarc[
			decoration = {
				text along path,
				text = {|\Huge\bfseries\sffamily\color{white}|MINIMAL},
				text align = center, raise=\scale* 0.225 cm
			}, 
			decorate
%			% for debug only
%			draw=black,
%			postaction={decorate}
		](0,0)(180:0:\scale*4.125cm);
		
		% curved text bottom
		\centerarc[
			decoration={text along path, text align = center,
				text={|\Huge\bfseries\sffamily\color{white}|THESIS}, raise = -0.75 cm%raise=0.125cm
			},
			decorate
%			% for debug only
%			draw=black,
%			postaction={decorate}
		](0,0)(-145:-35:\scale*4.125cm)
		
	\end{tikzpicture}
\end{document}